\chapter{HARDWARE BRING-UP}
\label{ch:hardware}

This chapter records every change that was needed to make each device work on
the BeagleBone AI. The Robotics Cape was designed for the BeagleBone Black and
Blue, whose AM335x processor has a different pin multiplexer, different GPIO
numbering and different peripheral addresses from the AM5729 on the
BeagleBone AI. Almost every problem in this chapter comes from that mismatch.

\section{OPERATING SYSTEM AND BOOT}

The board runs the 2020-04-06 TI image (Debian~10 Buster, kernel
\texttt{4.14.108-ti-r136}). Debian~10 is out of support, so
\texttt{/etc/apt/sources.list} points at \texttt{archive.debian.org} with the
security line commented out; without that, \texttt{apt} cannot install
anything (it was needed for \texttt{libatlas-base-dev}, which depthai uses).

The boot configuration is \file{system/uEnv.txt}, installed as
\texttt{/boot/uEnv.txt}:
\begin{lstlisting}
uname_r=4.14.108-ti-r136
dtb=am5729-beagleboneai-roboticscape.dtb
cmdline=coherent_pool=1M
\end{lstlisting}
The \texttt{dtb=} line makes U-Boot load the custom Robotics Cape device tree
instead of the stock one. Journald was made persistent
(\texttt{/var/log/journal}) on 1 October 2026 so that the cause of an
unexpected reboot survives the reboot.

\section{DEVICE TREE}
\label{sec:dt}

The device tree is built from \texttt{beagleboard/BeagleBoard-DeviceTrees}
(branch \texttt{v4.14.x-ti}, commit \texttt{4a9c0a6}) plus
\file{devicetree/roboticscape-bbai.patch}. The full source files are also
kept in \file{devicetree/}. Table~\ref{tab:dt} lists every change and why it
was made.

\begin{table}[htbp]
\centering
\caption{Device tree changes on top of upstream}
\label{tab:dt}
\small
\begin{tabularx}{\linewidth}{>{\raggedright\arraybackslash}p{1.9in}X}
\toprule
Change & Reason \\
\midrule
\texttt{\&hdmi} disabled & HDMI shares pins with cape signals. \\
\texttt{\&i2c4} enabled on P9.19/P9.20 (mode 7), shared balls set to GPIO
  input & Cape I$^2$C2 header. The BB-AI routes two SoC balls to each of these
  header pins, so the unused ball must be parked. \\
P9.13 muxed to \texttt{pr2\_pru0\_gpo15} (mode 13) & Motor 1 direction~B. The
  AM335x GPIO this pin used is not available on the AM5729 ball that reaches
  the header; the only usable function is a PRU output (\S\ref{sec:pru}). \\
\texttt{PIN\_INPUT\_PULLUP} on all 8 encoder pins (eQEP1A/B, eQEP2A/B,
  eQEP3A/B, PRU1 gpi16/18) & The A1230 outputs are open-drain and need a
  pull-up. Without it no encoder counted at all (2026-10-01). \\
\texttt{\&uart6} enabled & Bluetooth HCI for the BCM43455. It had been
  commented out in the base \texttt{.dts} in 2022 for an unrecorded reason;
  re-enabled and tested on 2026-10-01. \\
\bottomrule
\end{tabularx}
\end{table}

The patch also moves the \texttt{epwmss2}/\texttt{ehrpwm2} and UART nodes
around without changing what they enable, and adds a model string.

\paragraph{Build and install.} On the board (needs \texttt{sudo} and a reboot):
\begin{lstlisting}
cd ~/src/BeagleBoard-DeviceTrees && git apply .../roboticscape-bbai.patch
make src/arm/am5729-beagleboneai-roboticscape.dtb
sudo cp /boot/dtbs/4.14.108-ti-r136/am5729-beagleboneai-roboticscape.dtb ~/roboticscape.dtb.bak
sudo cp src/arm/am5729-beagleboneai-roboticscape.dtb /boot/dtbs/4.14.108-ti-r136/
sudo reboot
\end{lstlisting}
The backup in the home folder is the rollback path; the board's eMMC can also
be repaired from a microSD boot if the new tree does not boot.

\section{SERIAL PORTS}

The cape's UART header names, the AM5729 UART instances and the Linux device
names do not line up, which caused the lidar to be looked for on the wrong
port at first. Table~\ref{tab:uarts} gives the confirmed mapping.

\begin{table}[htbp]
\centering
\caption{UART assignments (confirmed on the board)}
\label{tab:uarts}
\small
\begin{tabular}{llll}
\toprule
Device & Cape header & Linux device & Settings \\
\midrule
RPLIDAR A1 & UART1 & \texttt{/dev/ttyS3} (old name \texttt{ttyO3}) & \SI{115200}{baud} \\
EM-506 GPS & GPS & \texttt{/dev/ttyS2} (SoC UART3) & \SI{4800}{baud} 8N1 \\
Bluetooth HCI & (on board) & \texttt{/dev/ttyS5} (SoC UART6) & \SI{3}{Mbaud}, RTS/CTS \\
Debug console & (on board) & \texttt{ttyS0} & \SI{115200}{baud}, PC COM3 \\
\bottomrule
\end{tabular}
\end{table}

\section{LIBROBOTCONTROL FORK}
\label{sec:rclib}

Motors, encoders and the IMU are driven through \texttt{librobotcontrol}
(used from Python via \texttt{rcpy}). The upstream library supports the
BeagleBone AI only partially. The fork
\texttt{jstumfoll2/librobotcontrol}, branch \texttt{v1.1\_AIfixes}, is a
submodule of this repository and carries two commits on top of upstream
\texttt{99a08db}: Jason's 2022 work (\texttt{d90e24b}) and the uncommitted
board changes recovered on 2026-10-01 (\texttt{593182b}). The functional
changes are:
\begin{itemize}
\item \textbf{GPIO numbering} (\texttt{motor.c}). All motor direction pins and
  the motor standby pin were remapped to AM5729 GPIO banks (for example MDIR1A
  is \texttt{gpio5\_0}, MOT\_STBY is \texttt{gpio6\_20}). The per-channel
  polarity table was changed to $\{+1,-1,+1,-1\}$.
\item \textbf{A virtual GPIO chip for the PRU} (\texttt{gpio.c}). Chip number
  11 is redirected to the character device \texttt{/dev/rpmsg\_pru32}. Setting
  a pin on chip 11 writes the ASCII character \texttt{'0'} or \texttt{'1'} to
  that device instead of issuing a GPIO ioctl. Motor 1's direction~B pin is
  defined as chip 11, so \texttt{rc\_motor\_init()} fails unless the PRU
  firmware below is running.
\item \textbf{eQEP sysfs paths} (\texttt{encoder\_eqep.c}). Updated to the
  AM5729 addresses (\texttt{4843e000.epwmss} and so on).
\item \textbf{IMU bus} (\texttt{mpu.c}). The MPU-9250 is on I$^2$C bus 3 with
  its interrupt on \texttt{gpio5\_17} (P9.25).
\item \textbf{PRU encoder} (\texttt{encoder\_pru.c}). Encoder channel 4 is
  counted by firmware on PRU1\_1 (\texttt{am57xx-pru1\_1-rc-encoder-fw-c}),
  read from PRU shared memory at offset 0.
\end{itemize}

\section{MOTOR 1 DIRECTION FIRMWARE ON PRU2\_0}
\label{sec:pru}

The PRU firmware \texttt{gpio\_rpmsg.pru2\_0.c}
(\file{librobotcontrol/pru\_firmware/customfw/}) is about 100 lines. It opens
an RPMsg channel named \texttt{rpmsg-pru} on port 32, which the kernel exposes
as \texttt{/dev/rpmsg\_pru32}. Each message received is parsed with
\texttt{atoi}; a non-zero value sets bit~15 of the PRU output register
\texttt{R30} (P9.13), and zero clears it. After each message the firmware
waits $\num{10200}+\num{4000000}$ PRU cycles, about \SI{20}{ms} at
\SI{200}{MHz}, so motor 1 can reverse direction at most about 50 times per
second. That is far faster than any command the robot sends.

The firmware is installed as \texttt{/lib/firmware/am57xx-pru2\_0-fw}. It is
not started at boot; \texttt{base\_node.py} starts it before initialising the
motors, equivalent to
\begin{lstlisting}
echo am57xx-pru2_0-fw > /sys/class/remoteproc/remoteproc2/firmware
echo start            > /sys/class/remoteproc/remoteproc2/state
\end{lstlisting}
and then waits up to \SI{5}{s} for \texttt{/dev/rpmsg\_pru32} to appear.

\section{WHEEL ENCODERS}

Each wheel has a DAGU 8-pole magnet ring (4 north, 4 south) read by an Allegro
A1230 dual-channel Hall sensor in quadrature. Two facts from the datasheets
explained why nothing counted at first:
\begin{itemize}
\item The A1230 outputs are open-drain, and the datasheet requires pull-up
  resistors. The cape has none on these lines, so the internal pull-ups were
  enabled in the device tree (Table~\ref{tab:dt}).
\item The A1230 needs at least \SI{3.3}{V}; the cape's encoder headers supply
  exactly \SI{3.3}{V}, which is in specification but at the bottom of the range.
\end{itemize}

An 8-pole ring with a quadrature sensor gives $4\times 4 = 16$ edges per
revolution in principle. A \SI{3}{ft} (\SI{0.914}{m}) push measured about 55
counts on every wheel, all four within 2 counts of each other, giving
\num{15.0} counts per revolution, and that value is in \file{base.yaml}.

\begin{table}[htbp]
\centering
\caption{Encoder and motor channel mapping (spin and lifted-wheel tests, 2026-10-01)}
\label{tab:channels}
\small
\begin{tabular}{lcccc}
\toprule
Wheel & Encoder channel & Counter & Encoder sign & Motor channel \\
\midrule
Left back   & 1 & eQEP & $+1$ & M1 \\
Right back  & 2 & eQEP & $-1$ & M4 \\
Left front  & 3 & eQEP & $+1$ & M2 \\
Right front & 4 & PRU1\_1 & $-1$ & M3 \\
\bottomrule
\end{tabular}
\end{table}

All four motors drive forward with positive duty (\texttt{motor\_signs} all
$+1$), after the library's own polarity table. The right-side encoders count
down when driving forward because they are mounted mirrored, which the
encoder signs correct.

\section{INERTIAL MEASUREMENT UNIT}

The MPU-9250 is read through \texttt{rcpy.mpu9250} with the magnetometer
enabled. \texttt{rcpy} returns acceleration in \si{m/s^2}, angular rate in
\si{\degree/s} and magnetic field in \si{\micro T}; \texttt{base\_node.py}
converts to SI (\si{rad/s}, \si{T}) before publishing. Relevant datasheet
values at the default \SI{\pm250}{\degree/s} range are a rate-noise spectral
density of \SI{0.01}{\degree/s/\sqrt{Hz}}, an initial zero-rate offset of up to
\SI{\pm5}{\degree/s}, and a zero-rate offset that can move by up to
\SI{\pm30}{\degree/s} over the full temperature range. The last figure is why
the gyroscope bias must be measured at every start (\S\ref{sec:gyro}).

The magnetometer read about \SI{150}{\micro T} on the robot, against Earth's
roughly \SI{50}{\micro T}. The extra field comes from the motor magnets,
the encoder magnet rings and current in the wiring, which is why it needs the
calibration in \S\ref{sec:mag}.

\section{LIDAR}

The RPLIDAR A1M8 is wired to the cape's UART1 header and appears as
\texttt{/dev/ttyS3}. It reports healthy and publishes \topic{/scan} at about
\SI{7.2}{Hz} through \texttt{rplidar\_ros} (catkin workspace
\texttt{\textasciitilde/melodic-catkin-ws}). The A1 scans clockwise; the
driver converts to the counter-clockwise ROS convention.

\section{GPS}

The EM-506 is a SiRF Star IV receiver that outputs NMEA 0183 at
\SI{4800}{baud} on \texttt{/dev/ttyS2}. Two pieces of software read it:
\begin{itemize}
\item \file{gps/libgps-bbai-ttyS2-4800.patch}: a patch to the small C library
  \texttt{wdalmut/libgps} that changes its port and baud rate. This was the
  2022 test path.
\item \file{ros/bbai\_base/scripts/gps\_node.py}: the ROS publisher used now
  (\S\ref{sec:gpsnode}). It configures the UART itself with \texttt{termios}.
\end{itemize}
Indoors the receiver reports no fix; a fix needs open sky.

\section{BLUETOOTH}
\label{sec:bt}

The on-board BCM43455 Bluetooth controller talks HCI over SoC UART6. Three
things were needed to bring it up:
\begin{enumerate}
\item \texttt{\&uart6} enabled in the device tree (Table~\ref{tab:dt}). It then
  appears as \texttt{/dev/ttyS5}.
\item \file{system/blacklist-btsdio.conf} in \texttt{/etc/modprobe.d/}. The
  \texttt{btsdio} driver otherwise creates a non-working \texttt{hci0} on the
  SDIO bus, so the real controller would become \texttt{hci1}.
\item \file{ros/bbai\_teleop/bt-up.sh}, run at boot by
  \texttt{bbai-bluetooth.service}. It finds the tty whose device is
  \texttt{48068000.serial}, attaches it with
  \texttt{hciattach <tty> bcm43xx 3000000 flow} (retrying at \SI{3}{Mbaud} if
  the chip was already initialised), brings \texttt{hci0} up, disables sniff
  mode on new links and turns off Wi-Fi power saving. The last two keep the
  controller's input reports at full rate, because the BCM43455 shares one
  radio between Wi-Fi and Bluetooth.
\end{enumerate}
The PS4 controller is paired once with \file{pair-ps4.sh}, which runs the
whole \texttt{bluetoothctl} conversation (agent, scan, pair, trust, connect)
in a single session. The kernel's \texttt{hid-sony} driver then creates
\texttt{/dev/input/js0}.

\section{OAK-D-LITE CAMERA}

The camera runs with depthai~2.19.1 on the board's Python~3.7 over USB~2, at
about \SI{15}{fps} for RGB plus depth. MobileNet-SSD runs on the camera's
Myriad X and detects the dog at 93--99\% confidence, about 10 detections per
second. The camera's EEPROM calibration was found blank; it was
self-calibrated on the PC with depthai~v2 tools and flashed on 2026-10-01
(epipolar error 0.36, baseline \SI{7.43}{cm}). The procedure is in
\file{docs/CALIBRATION.md} and the backup is in \file{camera/pc-tools/}.
Chapter~\ref{ch:camera} covers the camera's geometry, depth and
calibration in detail.

\section{POWER}
\label{sec:power}

The board reset several times on 1 October 2026. Persistent journald, a
once-per-second battery logger (\texttt{\textasciitilde/power.log}) and the
serial console narrowed the cause down:
\begin{itemize}
\item The battery input held steady (\SIrange{7.7}{8.2}{V}) through resets.
\item With the motor driver and lidar stopped, the camera ran detection for
  \SI{90}{s} with no reset.
\item The console showed the camera's USB disconnecting and reconnecting
  immediately before a reset, with no kernel panic.
\end{itemize}
The conclusion is that the cape's \SI{5}{V} rail, which feeds the BeagleBone
and its USB port, drops out when the camera (\SIrange{1}{1.5}{A}), motors,
lidar and Bluetooth all run together. The load estimate behind this is in
\S\ref{sec:powerbudget}. Later tests with a fast logger
(\S\ref{sec:resettests}) showed that the camera alone is enough to reset
the board, given a couple of minutes. A USB Y-cable did not help; the fix
was a USB-C hub on the BeagleBone's USB-C port with the battery bank on the
hub's power input.
